\chapter{Chép Bài Giải Không Hiểu}
\chapterintrobi{Trang vở rất đầy nhưng đầu vẫn còn trống}{Copy the solution without understanding}{Đây là kiểu học tạo cảm giác tiến bộ bằng lượng chữ, không bằng độ hiểu.}

\begin{lucbatbox}{Lục bát: Chữ nhiều mà nghĩa mỏng}
Bạn kia giải rất trơn tru\
Em chép theo hết từng khuông từng dòng\
Vở nhìn đầy đặn lạ không\
Hỏi lại bước nọ vẫn mông lung hoài\
Nếu tay chậm bớt một vài\
Đầu kịp hỏi lại, ngày mai đỡ mù
\end{lucbatbox}

\begin{tutuyetbox}{Thất ngôn tứ tuyệt: Chép đẹp chưa đủ}
Trang vở của em rất đẹp rồi\
Lời giải đi qua mượt lắm thôi\
Chỉ tiếc phần hiểu chưa kịp đi cùng chữ\
Nên gặp bài mới lại chơi vơi
\end{tutuyetbox}

\begin{ghichu}
Chép bài giải không hiểu là cách làm đầy vở nhanh. Nó không phải cách làm đầy tư duy nhanh tương ứng.
\end{ghichu}

\begin{englishnote}
\textbf{copy the solution}: chép lời giải.\par
\textbf{without understanding}: không hiểu thật.\par
\textbf{My notes are full, but I still do not get it.}: Vở em đầy mà em vẫn chưa hiểu.
\end{englishnote}